% --- 1. GEOMETRIA DELLA PAGINA ---
\usepackage{geometry}
\geometry{
    a4paper,
    total={170mm,257mm},
    left=15mm,
    right=15mm,
    top=20mm,
    bottom=20mm
}
\usepackage{float}

% --- 2. FONT E LEGGIBILITÀ ---
\usepackage[T1]{fontenc}
\usepackage[utf8]{inputenc}
\usepackage[italian]{babel}
\usepackage[sfdefault]{inter} 
\renewcommand{\familydefault}{\sfdefault}

% --- 3. DIAGRAMMI UML E GRAFICA ---
\usepackage{tikz}
\usepackage{pgf-umlsd} % Per i diagrammi di sequenza (SSD)
% NOTA: Se tikz-uml non è installato nel tuo sistema, questo pacchetto darà errore.
% Se lo usi, assicurati di avere il file tikz-uml.sty nella cartella del progetto.
\usepackage{tikz-uml} 

\usetikzlibrary{
    arrows.meta,
    shapes.symbols, % Fondamentale per la forma "note" (il foglietto)
    positioning,    % Per posizionare gli elementi "a destra di"
    calc
}

% DEFINIZIONE STILE NOTA (Così non devi scriverlo ogni volta)
\tikzset{
    note/.style={
        draw, 
        fill=yellow!10, % Un giallino leggero per farla sembrare una nota
        shape=note, 
        text width=4cm, 
        font=\scriptsize, 
        inner sep=7pt,
        shadow={opacity=0.2}
    }
}

% --- 4. CODICE SORGENTE ---
\usepackage{listings}
\usepackage{xcolor}
\definecolor{codegreen}{rgb}{0,0.6,0}
\lstset{
    basicstyle=\ttfamily\small,
    keywordstyle=\color{blue},
    stringstyle=\color{codegreen},
    breaklines=true,
    frame=single
}
\usepackage[normalem]{ulem}

% --- 5. ALTRI ESSENZIALI ---
\usepackage{amsmath, amssymb}
\usepackage{graphicx}
\graphicspath{ {./figures/} {../figures/} }
\usepackage{enumitem}
\usepackage{tabularx}
\setlist{nosep}

% Caricare hyperref SEMPRE per ultimo
\usepackage[colorlinks=true, linkcolor=blue]{hyperref}